% GENERATED FILE. Edit canonical lessons, then run npm run generate:books.
% canonical-source-hash: fnv1a64:84a7625499e60d4b
% canonical-lessons: KA-C267-gobbara, KA-C267-kunda, KA-C267-paip, KA-C267-pampu, KA-C267-getu

\chapter{Manure, Flowerpot, Pipe, Pump, Gate}
\label{ch:ka-manure-flowerpot-pipe-pump-gate}
\hlchaptermodality{267}

\begin{chapteropening}
\textbf{By the end of this chapter:} \emph{I can say manure, a flowerpot, a pipe, a pump and a gate.}

Complete the last lesson of chapter 267: I can say manure, a flowerpot, a pipe, a pump and a gate.
\end{chapteropening}

\section[\texorpdfstring{\kn{ಗೊಬ್ಬರ} (gobbara)}{ಗೊಬ್ಬರ (gobbara)}]{\kn{ಗೊಬ್ಬರ} (gobbara) — manure, fertiliser}

\label{lesson:KA-C267-gobbara}



\begin{quote}
Before the new one: say the Kannada for camphor, then the Kannada for copper.
\end{quote}

\subsection*{You'll want to know: \kn{ಗೊಬ್ಬರ}}
\textbf{\kn{ಗೊಬ್ಬರ}} — \emph{gobbara} — \textquotedblleft{}manure, fertiliser\textquotedblright{}.

\begin{grammarlens}[title={What you've built}]
One more word for this chapter.
\end{grammarlens}

\subsection*{Your turn}
\begin{itemize}
\raggedright
  \item \emph{Say it:} \emph{gobbara}
  \item \emph{Say it:} \emph{gobbara}, once more
  \item \emph{Say it:} say \emph{gobbara}
  \item \emph{Recall:} say the Kannada for camphor, then the Kannada for copper, then say \emph{gobbara} again
\end{itemize}

\begin{tcolorbox}[breakable,colback=teal!4,colframe=teal!35!black,title={Before you move on}]
What does \kn{ಗೊಬ್ಬರ} mean? (\textquotedblleft{}Manure, fertiliser\textquotedblright{}.) Say it once more.
\end{tcolorbox}
\section[\texorpdfstring{\kn{ಕುಂಡ} (kuṇḍa)}{ಕುಂಡ (kuṇḍa)}]{\kn{ಕುಂಡ} (kuṇḍa) — a flowerpot}

\label{lesson:KA-C267-kunda}



\begin{quote}
Before the new one: say the Kannada for copper, then the Kannada for manure.
\end{quote}

\subsection*{You'll want to know: \kn{ಕುಂಡ}}
\textbf{\kn{ಕುಂಡ}} — \emph{kuṇḍa} — \textquotedblleft{}a flowerpot\textquotedblright{}.

\begin{grammarlens}[title={What you've built}]
One more word for this chapter.
\end{grammarlens}

\subsection*{Your turn}
\begin{itemize}
\raggedright
  \item \emph{Say it:} \emph{kuṇḍa}
  \item \emph{Say it:} \emph{kuṇḍa}, once more
  \item \emph{Say it:} say \emph{kuṇḍa}
  \item \emph{Recall:} say the Kannada for copper, then the Kannada for manure, then say \emph{kuṇḍa} again
\end{itemize}

\begin{tcolorbox}[breakable,colback=teal!4,colframe=teal!35!black,title={Before you move on}]
What does \kn{ಕುಂಡ} mean? (\textquotedblleft{}A flowerpot\textquotedblright{}.) Say it once more.
\end{tcolorbox}
\section[\texorpdfstring{\kn{ಪೈಪ್} (paip)}{ಪೈಪ್ (paip)}]{\kn{ಪೈಪ್} (paip) — a pipe, a hose}

\label{lesson:KA-C267-paip}



\begin{quote}
Before the new one: say the Kannada for manure, then the Kannada for a flowerpot.
\end{quote}

\subsection*{You'll want to know: \kn{ಪೈಪ್}}
\textbf{\kn{ಪೈಪ್}} — \emph{paip} — \textquotedblleft{}a pipe, a hose\textquotedblright{}.

\begin{grammarlens}[title={What you've built}]
One more word for this chapter.
\end{grammarlens}

\subsection*{Your turn}
\begin{itemize}
\raggedright
  \item \emph{Say it:} \emph{paip}
  \item \emph{Say it:} \emph{paip}, once more
  \item \emph{Say it:} say \emph{paip}
  \item \emph{Recall:} say the Kannada for manure, then the Kannada for a flowerpot, then say \emph{paip} again
\end{itemize}

\begin{tcolorbox}[breakable,colback=teal!4,colframe=teal!35!black,title={Before you move on}]
What does \kn{ಪೈಪ್} mean? (\textquotedblleft{}A pipe, a hose\textquotedblright{}.) Say it once more.
\end{tcolorbox}
\section[\texorpdfstring{\kn{ಪಂಪು} (pampu)}{ಪಂಪು (pampu)}]{\kn{ಪಂಪು} (pampu) — a pump}

\label{lesson:KA-C267-pampu}



\begin{quote}
Before the new one: say the Kannada for a flowerpot, then the Kannada for a pipe.
\end{quote}

\subsection*{You'll want to know: \kn{ಪಂಪು}}
\textbf{\kn{ಪಂಪು}} — \emph{pampu} — \textquotedblleft{}a pump\textquotedblright{}.

\begin{grammarlens}[title={What you've built}]
One more word for this chapter.
\end{grammarlens}

\subsection*{Your turn}
\begin{itemize}
\raggedright
  \item \emph{Say it:} \emph{pampu}
  \item \emph{Say it:} \emph{pampu}, once more
  \item \emph{Say it:} say \emph{pampu}
  \item \emph{Recall:} say the Kannada for a flowerpot, then the Kannada for a pipe, then say \emph{pampu} again
\end{itemize}

\begin{tcolorbox}[breakable,colback=teal!4,colframe=teal!35!black,title={Before you move on}]
What does \kn{ಪಂಪು} mean? (\textquotedblleft{}A pump\textquotedblright{}.) Say it once more.
\end{tcolorbox}
\section[\texorpdfstring{\kn{ಗೇಟು} (gēṭu)}{ಗೇಟು (gēṭu)}]{\kn{ಗೇಟು} (gēṭu) — a gate}

\label{lesson:KA-C267-getu}



\begin{quote}
Before the new one: say the Kannada for a pipe, then the Kannada for a pump.
\end{quote}

\subsection*{You'll want to know: \kn{ಗೇಟು}}
\textbf{\kn{ಗೇಟು}} — \emph{gēṭu} — \textquotedblleft{}a gate\textquotedblright{}.

\begin{grammarlens}[title={What you've built}]
One more word for this chapter.
\end{grammarlens}

\subsection*{Your turn}
\begin{itemize}
\raggedright
  \item \emph{Say it:} \emph{gēṭu}
  \item \emph{Say it:} \emph{gēṭu}, once more
  \item \emph{Say it:} say \emph{gēṭu}
  \item \emph{Recall:} say the Kannada for a pipe, then the Kannada for a pump, then say \emph{gēṭu} again
\end{itemize}

\begin{tcolorbox}[breakable,colback=teal!4,colframe=teal!35!black,title={Before you move on}]
What does \kn{ಗೇಟು} mean? (\textquotedblleft{}A gate\textquotedblright{}.) Say it once more.
\end{tcolorbox}
